\begin{lemma}[A concentrated index-sum class]\label{tou:concentration}
Let $A=\binom mD$ and $V=D(m-D)(m+1)$.
The largest index-sum class of $D$-subsets of
$\{1,\ldots,m\}$ has size $J$ satisfying
\[
 J\ge \left\lceil\frac{A}{\sqrt{V+1}}\right\rceil.
\]
The same bound holds for any translated interval of indices.
In particular $J=\Omega(A/m^{3/2})$ at fixed density $D/m$.
\end{lemma}
\begin{proof}
For a uniform $D$-subset, its index sum $S$ has variance $V/12$.
This follows by writing $S$ as the sum of index-weighted selection
indicators: their means are $D/m$ and distinct pairs have joint mean
$D(D-1)/(m(m-1))$.
Let $U$ be independent and uniform on $[-1/2,1/2]$.
The density of $Y=S+U$ is at most $\mu=J/A$, and its variance is
$(V+1)/12$. Any real random variable with density at most $\mu$
has variance at least $1/(12\mu^2)$: at distance $t$ from its mean,
its tail probability is at least $1-2\mu t$, and hence
\[
 \operatorname{Var}(Y)\ge
 \int_0^{1/(2\mu)}2t(1-2\mu t)\,dt=\frac1{12\mu^2}.
\]
Consequently $J^2(V+1)\ge A^2$, proving the claim.
Translation changes all index sums by the same amount.
\end{proof}
